\chapter{Frontend Web com CloudFront para Operação da API e Ingestão RAG}
\label{chap:frontend_cloudfront_console}

% ---------------------------------------------------------------
\section{Escopo do cap\'itulo}
% ---------------------------------------------------------------

Este cap\'itulo apresenta a camada de experi\^encia do usu\'ario da
plataforma: uma interface web est\'atica distribu\'ida por
\textit{Amazon CloudFront}\index{CloudFront|textbf}\footnote{%
  Rede de distribui\c{c}\~ao de conte\'udo (CDN) da AWS utilizada para
  publicar arquivos est\'aticos com baixa lat\^encia e alta disponibilidade.}
com origem em bucket \textit{Amazon S3}\index{S3 (AWS)} privado,
destinada a dois fluxos principais: intera\c{c}\~ao com a API ass\'incrona
via \texttt{POST /query} e \texttt{GET /status/\{request\_id\}} e ingest\~ao
controlada de documentos no pipeline \textit{Retrieval-Augmented Generation}
(RAG\index{RAG|textbf}) via \texttt{POST /documentos}. O objetivo \'e
fornecer uma superf\'icie web de opera\c{c}\~ao que possa ser utilizada por
usu\'arios de neg\'ocio, analistas de homologa\c{c}\~ao e equipes t\'ecnicas sem
necessidade de notebooks ou scripts auxiliares.

% ---------------------------------------------------------------
\section{Objetivos de aprendizagem}
% ---------------------------------------------------------------

Ao concluir este cap\'itulo, o leitor ser\'a capaz de:

\begin{itemize}[leftmargin=*]
  \item Descrever a arquitetura de publica\c{c}\~ao de um frontend
        est\'atico via CloudFront e bucket S3 privado com
        \textit{Origin Access Control}\index{Origin Access Control|textbf};
  \item Implementar uma interface web para envio de consultas ao agente
        e acompanhamento do processamento ass\'incrono;
  \item Disponibilizar, na mesma interface, um fluxo de envio de
        documentos textuais para indexa\c{c}\~ao no RAG;
  \item Avaliar limita\c{c}\~oes do backend atual para formatos bin\'arios
        como PDF e DOCX e propor a extens\~ao arquitetural adequada.
\end{itemize}

% ---------------------------------------------------------------
\section{Conceitos te\'oricos}
% ---------------------------------------------------------------

\subsection{Frontend est\'atico como superf\'icie de opera\c{c}\~ao}

Em arquiteturas banc\'arias modernas, a separa\c{c}\~ao entre camada de
apresenta\c{c}\~ao e servi\c{c}os de neg\'ocio reduz acoplamento, melhora a
escalabilidade e simplifica a estrat\'egia de publica\c{c}\~ao. Para o caso
em estudo, um frontend est\'atico atende adequadamente aos requisitos
porque:

\begin{itemize}[leftmargin=*]
  \item o processamento de neg\'ocio ocorre integralmente no backend
        ass\'incrono (API Gateway + Lambda + SQS + DynamoDB);
  \item a UI atua apenas como orquestradora de chamadas HTTP,
        renderiza\c{c}\~ao de respostas e coleta de insumos do usu\'ario;
  \item a camada est\'atica pode ser publicada sem servidor dedicado,
        com custo operacional reduzido e baixa superf\'icie de ataque.
\end{itemize}

\subsection{CloudFront + S3 privado}

A distribui\c{c}\~ao CloudFront \'e configurada para apontar para um bucket
S3 privado que armazena os arquivos \texttt{index.html}, \texttt{styles.css},
\texttt{app.js} e \texttt{config.js}. O acesso direto ao bucket permanece
bloqueado; apenas a distribui\c{c}\~ao CloudFront, por meio de
\textit{Origin Access Control} (OAC), possui permiss\~ao de leitura.
Esse padr\~ao \'e prefer\'ivel ao uso de \textit{website endpoint}
p\'ublico do S3, pois preserva o princ\'ipio de exposi\c{c}\~ao m\'inima e
centraliza a entrega HTTPS.

\subsection{Chat ass\'incrono e ingest\~ao RAG em uma \~unica UI}

A interface implementa dois fluxos distintos:

\begin{enumerate}[leftmargin=*]
  \item \textbf{Consulta ass\'incrona:} o usu\'ario preenche pergunta,
        modo de opera\c{c}\~ao e, opcionalmente, dados do cliente; a UI
        envia o payload para \texttt{/query}, recebe um
        \texttt{request\_id}\index{request\_id} e executa \textit{polling}
        em \texttt{/status/\{request\_id\}} at\'e o estado \texttt{COMPLETED}.
  \item \textbf{Ingest\~ao de documentos:} o usu\'ario informa o tipo
        (\texttt{documento} ou \texttt{entrevista}), t\'itulo e conte\'udo,
        podendo carregar um arquivo textual local. O frontend converte o
        arquivo em texto e envia JSON para \texttt{/documentos}.
\end{enumerate}

A limita\c{c}\~ao atual decorre do contrato do endpoint \texttt{/documentos},
que aceita apenas JSON com campo \texttt{texto}; n\~ao h\'a suporte nativo a
\textit{multipart/form-data}, nem extra\c{c}\~ao de conte\'udo de formatos
bin\'arios como PDF. Assim, a UI suporta apenas arquivos textuais e colagem
manual de conte\'udo.

% ---------------------------------------------------------------
\section{Implementa\c{c}\~ao orientada por passos}
% ---------------------------------------------------------------

\subsection{Passo 1 --- Estruturar os assets do frontend}

A aplica\c{c}\~ao web foi implementada sem \textit{framework} JavaScript,
com o objetivo de minimizar depend\^encias e facilitar a publica\c{c}\~ao:

\begin{lstlisting}[language=bash,
  caption={Estrutura dos arquivos est\'aticos do console web.},
  label={lst:frontend_assets}]
frontend/
  cloudfront_console/
    index.html
    styles.css
    app.js
    config.js
\end{lstlisting}

O arquivo \texttt{config.js}\index{config.js|textbf} \'e sobrescrito no deploy
Terraform com a URL real da API Gateway, eliminando a necessidade de
compila\c{c}\~ao por ambiente.

\subsection{Passo 2 --- Publicar a UI em CloudFront}

O m\'odulo Terraform \texttt{cloudfront\_frontend} cria:

\begin{itemize}[leftmargin=*]
  \item bucket S3 privado com versionamento e criptografia;
  \item OAC para leitura segura do bucket pela distribui\c{c}\~ao;
  \item distribui\c{c}\~ao CloudFront com \texttt{default\_root\_object = index.html};
  \item objetos S3 contendo os assets do frontend e o \texttt{config.js}
        parametrizado.
\end{itemize}

\begin{lstlisting}[language=bash,
  caption={Trecho do m\'odulo Terraform do frontend CloudFront.},
  label={lst:terraform_cloudfront_frontend}]
module "cloudfront_frontend" {
  source       = "./modules/cloudfront_frontend"
  project_name = var.project_name
  environment  = var.environment
  api_endpoint = module.api_gateway.api_endpoint
}

output "frontend_url" {
  value = "https://${module.cloudfront_frontend.distribution_domain_name}"
}
\end{lstlisting}

\subsection{Passo 3 --- Implementar o fluxo de consulta ass\'incrona}

A interface web coleta pergunta, modo e dados do cliente e monta o payload
compat\'ivel com o contrato OpenAPI da aplica\c{c}\~ao. Em seguida:

\begin{enumerate}[leftmargin=*]
  \item envia \texttt{POST /query};
  \item armazena o \texttt{request\_id};
  \item inicia \textit{polling} com \textit{backoff} exponencial;
  \item atualiza a UI com status intermedi\'arios e resultado final.
\end{enumerate}

Esse padr\~ao desacopla a experi\^encia do usu\'ario da dura\c{c}\~ao real do
pipeline RAG, compatibilizando a interface web com o limite de 29 segundos
do API Gateway descrito no Cap\'itulo~\ref{chap:pipeline_assincrono_aws}.

\subsection{Passo 4 --- Implementar envio de documentos para RAG}

O fluxo de ingest\~ao aceita dois modos de entrada:

\begin{itemize}[leftmargin=*]
  \item \textbf{Upload de arquivo textual:} a UI l\^e o arquivo no navegador
        com \texttt{File.text()} e preenche o campo \texttt{texto};
  \item \textbf{Colagem manual:} o usu\'ario insere diretamente o conte\'udo
        no \textit{textarea}, permitindo testes e indexa\c{c}\~ao r\'apida.
\end{itemize}

O payload enviado para a API segue o formato:

\begin{lstlisting}[language=bash,
  caption={Payload de ingest\~ao de documento no RAG.},
  label={lst:payload_documento_rag}]
{
  "tipo": "documento",
  "titulo": "Politica de Investimentos 2026",
  "texto": "Clientes conservadores priorizam liquidez..."
}
\end{lstlisting}

Quando o tipo \'e \texttt{entrevista}, a UI pode enviar adicionalmente
\texttt{cliente\_id} ou \texttt{cluster\_id}, alinhando a indexa\c{c}\~ao com
os modos \textit{twin}, \textit{persona} e \textit{segmento}.

\subsection{Passo 5 --- Planejar evolu\c{c}\~ao para PDF e DOCX}

Para suportar documentos bin\'arios, recomenda-se introduzir uma etapa de
extra\c{c}\~ao antes da chamada a \texttt{/documentos}. Duas abordagens s\~ao
plaus\'iveis:

\begin{enumerate}[leftmargin=*]
  \item \textbf{Extra\c{c}\~ao no backend:} novo endpoint de upload recebe o
        bin\'ario, armazena em S3 e aciona Lambda com parser (Apache Tika,
        PyMuPDF, python-docx ou Amazon Textract\index{Textract|textbf}).
  \item \textbf{Pr\'e-processamento cliente:} alternativa limitada a alguns
        formatos e menos indicada para ambiente banc\'ario, pois exp\~oe a
        extra\c{c}\~ao ao navegador.
\end{enumerate}

Para o cen\'ario banc\'ario, a primeira abordagem \'e superior por motivos
de rastreabilidade, auditoria e controle de seguran\c{c}a.

% ---------------------------------------------------------------
\section{Plano de testes do cap\'itulo}
% ---------------------------------------------------------------

\subsection{Teste funcional}

\begin{itemize}[leftmargin=*]
  \item \textbf{TC-01:} Abrir a URL do CloudFront e verificar o
        carregamento de \texttt{index.html}, \texttt{styles.css} e
        \texttt{config.js} sem acesso direto p\'ublico ao bucket S3.
  \item \textbf{TC-02:} Executar uma consulta em modo \texttt{segmento}
        e verificar que a UI exibe \texttt{request\_id}, status intermedi\'ario
        e resposta final.
  \item \textbf{TC-03:} Submeter um documento textual e verificar resposta
        HTTP 200 com \texttt{chunks\_indexados > 0}.
  \item \textbf{TC-04:} Tentar enviar um arquivo PDF e verificar que a UI
        informa explicitamente a limita\c{c}\~ao atual do backend.
\end{itemize}

\subsection{Teste de valida\c{c}\~ao em produ\c{c}\~ao (AWS)}

Foi executado um smoke test completo no ambiente publicado, com os
seguintes endere\c{c}os:

\begin{itemize}[leftmargin=*]
  \item \textbf{Frontend CloudFront:}
        \texttt{https://djxkmivr2akfz.cloudfront.net};
  \item \textbf{API Gateway:}
        \texttt{https://3vqrahlmyg.execute-api.sa-east-1.amazonaws.com}.
\end{itemize}

Os resultados observados foram:

\begin{itemize}[leftmargin=*]
  \item \textbf{TVP-01 (ingest\~ao entrevista):}
        \texttt{ok=true}, \texttt{chunks\_indexados=1};
  \item \textbf{TVP-02 (consulta persona):}
        \texttt{request\_id}: \nolinkurl{ddbc5210-63fc-47bc-a274-283e60be84bc}. \par
        Estado final \texttt{COMPLETED};
  \item \textbf{TVP-03 (consulta twin):}
        \texttt{request\_id}: \nolinkurl{e5e348e7-d266-4ae9-9f26-eb0e659a56a7}. \par
        Estado final \texttt{COMPLETED};
  \item \textbf{TVP-04 (corre\c{c}\~ao de exibi\c{c}\~ao da resposta):}
        Identificado e corrigido defeito no \texttt{app.js} em que o campo
        \texttt{data.result} era exibido diretamente como texto, resultando
        em \texttt{[object Object]}. A corre\c{c}\~ao extrai
        \texttt{result.resposta} ap\'os \textit{parse} JSON, exibindo o
        texto leg\'ivel do agente no painel \textbf{Resposta} e os metadados
        (\texttt{cliente\_id}, \texttt{cluster\_id}, \texttt{confiabilidade})
        no painel \textbf{Status}. Implantado em produ\c{c}\~ao via
        \texttt{aws s3 sync} e invalida\c{c}\~ao do cache CloudFront em
        02/05/2026.
\end{itemize}

Com isso, verifica-se a opera\c{c}\~ao ponta a ponta da interface web em
produ\c{c}\~ao, incluindo ingest\~ao RAG e consulta ass\'incrona com
\textit{polling} at\'e a conclus\~ao.

\subsection{Teste de seguran\c{c}a}

\begin{itemize}[leftmargin=*]
  \item \textbf{TS-01:} Confirmar que o bucket S3 do frontend possui
        \textit{public access block} habilitado.
  \item \textbf{TS-02:} Verificar que apenas a distribui\c{c}\~ao CloudFront
        possui permiss\~ao de leitura nos objetos est\'aticos.
  \item \textbf{TS-03:} Validar cabe\c{c}alhos CORS da API para
        permitir consumo pelo frontend, sem abrir al\'em do necess\'ario.
\end{itemize}

% ---------------------------------------------------------------
\section{Crit\'erios de aceite}
% ---------------------------------------------------------------

\begin{itemize}[leftmargin=*]
  \item TC-01 a TC-04 aprovados.
  \item TS-01 a TS-03 aprovados.
  \item URL \texttt{frontend\_url} gerada pelo Terraform e acess\'ivel via HTTPS.
  \item Fluxos de consulta e ingest\~ao operacionais em \texttt{dev}.
  \item Smoke test em produ\c{c}\~ao aprovado para ingest\~ao,
        \texttt{persona} e \texttt{twin} (TVP-01 a TVP-03).
  \item Resposta leg\'ivel do agente exibida no painel \textbf{Resposta}
        via \texttt{result.resposta} (TVP-04 --- defei\-to de exibi\c{c}\~ao corrigido).
\end{itemize}

% ---------------------------------------------------------------
\section{Espa\c{c}o para expans\~ao iterativa}
% ---------------------------------------------------------------

\begin{itemize}[leftmargin=*]
  \item \textbf{v17.1} --- Integrar autentica\c{c}\~ao com Amazon Cognito
        \index{Cognito|textbf} para restringir o acesso \`a interface web.
  \item \textbf{v17.2} --- Adicionar hist\'orico de consultas e visualiza\c{c}\~ao
        de status por usu\'ario autenticado.
  \item \textbf{v17.3} --- Suportar upload bin\'ario com pipeline de extra\c{c}\~ao
        para PDF, DOCX e imagens escaneadas.
  \item \textbf{Lacuna} --- Avaliar \textit{WebSocket API} como alternativa
        ao \textit{polling} HTTP para feedback de progresso em tempo real.
\end{itemize}
